\subsection{NILPOTENT RADICAL OF A LIE ALGEBRA}

DEFINITION 3. Let g be a Lie algebra. The nilpotent radical of g is the intersection of the kernels of the finite-dimensional simple representations of g.

Remarks. (1) Let \(\mathfrak{s}\) be the nilpotent radical of \(\mathfrak{g}\) . As every decreasing sequence of vector subspaces of \(\mathfrak{g}\) is stationary, there exists a finite number of finite-dimensional simple representations of \(\mathfrak{g}\) whose kernels have intersection \(\mathfrak{s}\) .

The direct sum of these representations is semi-simple and has kernel s. It follows that the set of kernels of finite-dimensional semi-simple representations of g has a least element, namely s.

\begin{enumerate}
\def\labelenumi{(\arabic{enumi})}
\setcounter{enumi}{1}
\item
  By Proposition 4 (c) of §4, no. 3, s is also the intersection of the largest nilpotency ideals of the finite-dimensional representations of g. In particular, s is contained in the largest nilpotent ideal of g and is therefore a nilpotent ideal of g.
\item
  Every linear form \(\lambda\) on \(\mathfrak{g}\) which is zero on \(\mathcal{D}\mathfrak{g}\) is a simple representation (with space \(\mathbf{K}\) ) of \(\mathfrak{g}\) , whence \(\lambda (\mathfrak{s}) = \{0\}\) . It follows that \(\mathfrak{s} \subset \mathcal{D}\mathfrak{g}\) . On the other hand, \(\mathfrak{s}\) is contained in the radical \(\mathfrak{r}\) of \(\mathfrak{g}\) by Remark 2. We shall prove that \(\mathfrak{s} = \mathfrak{r} \cap \mathcal{D}\mathfrak{g}\) .
\end{enumerate}

Lemma 1. Let \(\mathbf{V}\) be a finite-dimensional vector space over \(\mathbf{K}\) , \(\mathfrak{g}\) a subalgebra of \(\mathfrak{gl}(\mathbf{V})\) such that \(\mathbf{V}\) is a simple \(\mathfrak{g}\) -module and \(\mathfrak{a}\) a commutative ideal of \(\mathfrak{g}\) . Then \(\mathfrak{a} \cap \mathcal{D}\mathfrak{g} = \{0\}\) . Let \(\mathbf{S}\) be the subalgebra of \(\mathcal{L}(\mathbf{V})\) generated by 1 and \(\mathfrak{a}\) .

If \(\mathfrak{b}\) is an ideal of \(\mathfrak{g}\) contained in \(\mathfrak{a}\) such that \(\operatorname{Tr}bs = 0\) for all \(b\in \mathfrak{b}\) and all \(s\in S\) , then in particular, by definition of \(S\) , \(\operatorname{Tr}(b^n) = 0\) for every integer \(n > 0\) and hence \(b\) is nilpotent (Algebra, Chapter VII, § 5, no. 5, Corollary 4 to Proposition 13); as the elements of \(\mathfrak{b}\) are all nilpotent, \(\mathfrak{b} = \{0\}\) (§ 4, no. 3, Lemma 2). We first apply this to the ideal \([\mathfrak{g},\mathfrak{a}]\) of \(\mathfrak{g}\) . If \(x\in \mathfrak{g}\) , \(a\in \mathfrak{a}\) , \(s\in S\) , then \(\operatorname{Tr}[x,a]s = \operatorname{Tr}(xas - axs) = \operatorname{Tr}x(as - sa) = 0\) since \(as = sa\) ; hence \([\mathfrak{g},\mathfrak{a}] = \{0\}\) . Hence the elements of \(\mathfrak{g}\) commute with those of \(\mathfrak{a}\) and hence also with those of \(S\) . If \(x,y\) belong to \(\mathfrak{g}\) and \(s\in S\) , then

\[
\operatorname{Tr} [ x, y ] s = \operatorname{Tr} (x y s - y x s) = \operatorname{Tr} x (y s - s y) = 0
\]

since \(ys = sy\) ; then taking \(b\) to be the ideal \(\mathcal{D}\mathfrak{g} \cap a\) , it follows that \(\mathcal{D}\mathfrak{g} \cap a = \{0\}\) .

THEOREM 1. Let \(\mathfrak{g}\) be a Lie algebra, \(\mathfrak{r}\) its radical and \(\mathfrak{s}\) its nilpotent radical. Then \(\mathfrak{s} = \mathcal{D}\mathfrak{g} \cap \mathfrak{r}\) .

We already know that \(\mathfrak{s} \subset \mathcal{D}\mathfrak{g} \cap \mathfrak{r}\) . Hence it will suffice to show that if \(\rho\) is a finite-dimensional simple representation of \(\mathfrak{g}\) then \(\rho(\mathcal{D}\mathfrak{g} \cap \mathfrak{r}) = \{0\}\) . Let \(k\) be the least integer \(\geqslant 0\) such that \(\rho(\mathcal{D}^{k}\mathfrak{r}) = \{0\}\) ; we write \(\mathfrak{g}' = \rho(\mathfrak{g})\) , \(\mathfrak{a}' = \rho(\mathrm{D}^k\mathfrak{r})\) ; as \(\mathcal{D}^k\mathfrak{r}\) is an ideal of \(\mathfrak{g}\) , \(\mathfrak{a}'\) is an ideal of \(\mathfrak{g}'\) ; this ideal is commutative since \(\rho(\mathcal{D}^{k}\mathfrak{r}) = \{0\}\) . If \(\mathrm{V}\) is the space of \(\rho\) , \(\mathfrak{g}' \subset \mathfrak{gl}(\mathrm{V})\) and \(\mathrm{V}\) is a simple \(\mathfrak{g}'\) -module. Then \(\rho(\mathcal{D}\mathfrak{g} \cap \mathcal{D}^k\mathfrak{r}) \subset \mathcal{D}\mathfrak{g}' \cap \mathfrak{a}' = \{0\}\) . If \(k > 0\) , then \(\mathcal{D}^k\mathfrak{r} \subset \mathcal{D}\mathfrak{g}\) and \(\rho(\mathcal{D}^k\mathfrak{r}) = \{0\}\) , contrary to the definition of \(k\) . Hence \(k = 0\) , that is

\[
\rho (\mathcal{D}\mathfrak{g} \cap \mathfrak{r}) = \{0 \}.
\]

COROLLARY 1. Let g be a solvable Lie algebra. The nilpotent radical of g is \(\mathcal{D}\mathfrak{g}\) . If \(\rho\) is a finite-dimensional simple representation of g, \(\rho(\mathfrak{g})\) is commutative and the associative algebra L generated by 1 and \(\rho(\mathfrak{g})\) is a field of finite degree over K.

Here \(\mathfrak{r} = \mathfrak{g}\) , whence \(\mathfrak{s} = \mathcal{D}\mathfrak{g}\) . Hence \(\rho(\mathcal{D}\mathfrak{g}) = \{0\}\) , which shows that \(\mathfrak{g}' = \rho(\mathfrak{g})\) is commutative. Every element \(\neq 0\) of \(\mathbf{L}\) is invertible by Schur's Lemma; hence \(\mathbf{L}\) is a field.

COROLLARY 2 (Lie's Theorem). Let \(\mathfrak{g}\) be a solvable Lie algebra; suppose that \(\mathbf{K}\) is algebraically closed. Let \(\mathbf{M}\) be a \(\mathfrak{g}\) -module of finite dimension over \(\mathbf{K}\) and let \((\mathrm{M}_1)_{0 \leqslant i \leqslant r}\) be a Jordan-Hölder series of \(\mathbf{M}\) . Then \(\mathrm{M}_{i - 1} / \mathrm{M}_i\) is of dimension 1 over \(\mathbf{K}\) for \(1 \leqslant i \leqslant r\) and, for all \(x \in \mathfrak{g}\) , \(x_{\mathrm{M}_{i - 1} / \mathrm{M}_i} = \lambda_i(x)\) . 1, where \(\lambda_i\) is a linear form on \(\mathfrak{g}\) which is zero on \(\mathcal{D}\mathfrak{g}\) . In particular, every simple \(\mathfrak{g}\) -module of finite dimension over \(\mathbf{K}\) is in fact of dimension 1.

Let \(\rho_{i}\) be the representation of \(g\) on \(\mathbf{M}_{i-1}/\mathbf{M}_{i}\) . The associative algebra \(\mathbf{L}_{i}\) generated by 1 and \(\rho_{i}(g)\) is a field, a finite extension of K and therefore equal to K; and \(\mathbf{M}_{i-1}/\mathbf{M}_{i}\) is a simple \(\mathbf{L}_{i}\) -module, whence \(\dim \mathbf{M}_{i-1}/\mathbf{M}_{i} = 1\) . The rest of the corollary is obvious.

Remarks. (1) If \((\mathrm{M}_i)_{0\leqslant i\leqslant r}\) is replaced by another Jordan-Hölder series of M, the sequence \((\lambda_1,\dots ,\lambda_r)\) is replaced by a sequence of the form \((\lambda_{\pi (1)},\dots ,\lambda_{\pi (r)})\) where \(\pi\) is a permutation of \(\{1,\dots ,r\}\) , as follows from the Jordan-Hölder Theorem.

\begin{enumerate}
\def\labelenumi{(\arabic{enumi})}
\setcounter{enumi}{1}
\tightlist
\item
  Let \((e_1, \ldots, e_r)\) be a basis of M such that \(e_i \in \mathbf{M}_{i-1}, e_i \notin \mathbf{M}_i\) ( \(1 \leqslant i \leqslant r\) ). If \(x \in \mathfrak{g}\) the endomorphism of M corresponding to \(x\) is represented with respect to this basis by a triangular matrix whose diagonal coefficients are
\end{enumerate}

\[
\lambda_ {1} (x), \dots , \lambda_ {\tau} (x).
\]

COROLLARY 3. Suppose that K is algebraically closed. If g is an r-dimensional solvable Lie algebra, every ideal of g is a term of a decreasing sequence of ideals of dimensions r, r - 1, \ldots, 0.

Every ideal is part of a Jordan-Hölder series of g, considered as the space of the adjoint representation (Algebra, Chapter I, § 6, no. 14, Corollary to Theorem 8); then it suffices to apply Corollary 2.

COROLLARY 4. Suppose that \(\mathbf{K} = \mathbf{R}\) . Let \(\mathfrak{g}\) be a solvable Lie algebra. Every simple representation of \(\mathfrak{g}\) is of dimension \(\leqslant 2\) . Every ideal of \(\mathfrak{g}\) is a term of a decreasing sequence \((\mathfrak{g}_i)_{0\leqslant i\leqslant m}\) of ideals such that \(\mathfrak{g}_0 = \mathfrak{g},\mathfrak{g}_m = \{0\}\) , \(\dim \mathfrak{g}_{i - 1} / \mathfrak{g}_i\leqslant 2\) ( \(1\leqslant i\leqslant m\) ).

This is proved in a similar way to that for Corollaries 2 and 3, using the fact that every algebraic extension of \(\mathbf{R}\) is of degree \(\leqslant 2\) .

COROLLARY 5. For a Lie algebra g to be solvable, it is necessary and sufficient that \(\mathcal{D}\mathfrak{g}\) be nilpotent.

The condition is necessary by Corollary 1. It is sufficient since \(\mathfrak{g} / \mathcal{D}\mathfrak{g}\) is commutative.

COROLLARY 6. Let \(\rho\) be a finite-dimensional representation of a Lie algebra \(\mathfrak{g}\) . Let \(\mathfrak{r}\) be the radical of \(\mathfrak{g}\) . Every element \(x \in \mathfrak{r}\) such that \(\rho(x)\) is nilpotent belongs to the largest nilpotency ideal \(\mathfrak{n}\) of \(\rho\) .

Let \(V\) be the space of \(\rho\) ; let \(\langle V_i \rangle_{0 \leqslant i \leqslant r}\) be a Jordan-Hölder series for the \(r\) -module structure on \(V\) and let \(\rho_i\) be the representation of \(r\) with space \(V_i / V_{i-1}\) ( \(1 \leqslant i \leqslant r\) ). If \(\rho(x)\) is nilpotent, so is \(\rho_i(x)\) ; as for all \(i\) the algebra generated by \(\rho_i(x)\) is a field, \(\rho_i(x) = 0\) . Conversely, if \(\rho_i(x) = 0\) for all \(i, \rho(x) = 0\) . This shows that the set \(a\) of \(x \in r\) such that \(\rho(x)\) is nilpotent is an ideal of \(r\) . On the other hand, \([g, a] \subset \mathcal{D}g \cap r \subset n \cap r \subset a\) and hence \(a\) is an ideal of \(g\) . This proves that \(a \subset n\) .

COROLLARY 7. Let g be a Lie algebra and r its radical. The following four sets are identical: (a) the largest nilpotent ideal of g; (b) the largest nilpotent ideal of r; (c) the set of \(x \in \mathfrak{r}\) such that \(\operatorname{ad}_{\mathfrak{g}} x\) is nilpotent; (d) the set of \(x \in \mathfrak{r}\) such that \(\operatorname{ad}_{\mathfrak{r}} x\) is nilpotent.

Let these sets be denoted by \(\mathfrak{a}, \mathfrak{b}, \mathfrak{c}, \mathfrak{d}\) . The inclusions \(\mathfrak{a} \subset \mathfrak{b} \subset \mathfrak{d} \subset \mathfrak{c}\) are clear. \(\mathfrak{c} \subset \mathfrak{a}\) by Corollary 6 applied to the adjoint representation of \(\mathfrak{g}\) .

